\documentclass[11pt]{article}

\usepackage[british]{babel}
\usepackage[utf8]{inputenc}
\usepackage[T1]{fontenc}
\usepackage{lmodern}
\usepackage[margin=1in]{geometry}
\usepackage{amsmath,amssymb,amsthm}
\usepackage[round]{natbib}
\usepackage[hidelinks]{hyperref}

% Theorem environments, numbered within section with a shared counter.
\theoremstyle{plain}
\newtheorem{theorem}{Theorem}[section]
\newtheorem{lemma}[theorem]{Lemma}
\newtheorem{proposition}[theorem]{Proposition}
\newtheorem{corollary}[theorem]{Corollary}
\theoremstyle{definition}
\newtheorem{definition}[theorem]{Definition}
\newtheorem{assumption}[theorem]{Assumption}

\title{When do data-generated rubrics beat handcrafted prompts?}
\author{}
\date{\today}

\begin{document}
\maketitle

\section{Introduction}\label{sec:intro}

Language models are increasingly used as subroutines inside learning agents.
\citet{mitchell2026talk} describes such agents, which write rubrics from training data and then use them to label new inputs.
\citet{balcan2026talk} and \citet{balcan2025verifiers} develop a PAC framework for learning verifiers.
We combine the two views.
We treat a rubric as a hypothesis, the language model as a fixed map from rubrics to classifiers, and the generation of rubrics from data as a learning algorithm.
Our goal is to show when data-generated rubrics provably beat handcrafted prompts.

\section{Setup}\label{sec:setup}

Let $\mathcal X$ be an input space and $\mathcal Y=\{0,1\}$.
A task is a distribution $P$ on $\mathcal X\times\mathcal Y$.
Let $\mathcal R\subset\Sigma^*$ be a countable set of strings over a finite alphabet $\Sigma$, called prompts or rubrics.

\begin{definition}[LLM classifier]\label{def:llm-classifier}
A fixed language model $M$ assigns to each $r\in\mathcal R$ and $x\in\mathcal X$ a distribution $M(\cdot\mid r,x)$ on $\mathcal Y$.
The loss is $\ell(r,z)=1-M(y\mid r,x)\in[0,1]$ for $z=(x,y)$.
The risk is $R_P(r)=\mathbb E_{Z\sim P}\,\ell(r,Z)$.
The empirical risk on $S=(Z_1,\dots,Z_n)$ is $\hat R_S(r)=n^{-1}\sum_i\ell(r,Z_i)$.
The sampled indicator $\mathbb 1\{\hat Y_i\neq Y_i\}$ with $\hat Y_i\sim M(\cdot\mid r,X_i)$ has the same mean and range, and all bounds hold for it.
\end{definition}
% TODO(check): ollama_client.py obtains M(.|r,x) from the top-20 log-probabilities of the first generated token.
% It merges variants of "0" and "1" with a leading space, renormalises over {0,1}, and gives a label absent from the top 20 the smallest returned log-probability minus log 2.
% This is a valid M in the sense of the definition, but M then depends on the top-20 truncation.

\begin{definition}[Data-independent prompt]\label{def:data-independent}
$r_H\in\mathcal R$ is data-independent if it is fixed, or random and independent of the sample.
It may depend on a task description shared across tasks.
\end{definition}

\begin{definition}[Rubric procedure]\label{def:rubric-procedure}
Split an i.i.d.\ sample from $P$ into $S_1$ (size $n_1$) and $S_2$ (size $m$).
A proposal kernel $Q$ maps $S_1$ to a distribution $Q(\cdot\mid S_1)$ on $\mathcal R$.
Conditionally on $S_1$ draw $r_1,\dots,r_k$ i.i.d.\ from $Q(\cdot\mid S_1)$ and return $\hat r\in\arg\min_j\hat R_{S_2}(r_j)$.
Write $\varepsilon_m(k,\delta)=\sqrt{2\log(2k/\delta)/m}$.
\end{definition}

\section{Selection}\label{sec:selection}

\begin{proposition}[Selection]\label{prop:selection}
With probability at least $1-\delta$, $R_P(\hat r)\le\min_j R_P(r_j)+\varepsilon_m(k,\delta)$.
\end{proposition}

\begin{corollary}[Human prompt in the pool]\label{cor:human-in-pool}
If $r_H$ is added to the candidates, $R_P(\hat r)\le R_P(r_H)+\varepsilon_m(k+1,\delta)$ with probability at least $1-\delta$.
\end{corollary}

\begin{proposition}[Competitiveness through proposal mass]\label{prop:competitiveness}
Fix $\eta\ge0$ and $A_\eta=\{r:R_P(r)\le R_P(r_H)+\eta\}$.
Assume $Q(A_\eta\mid S_1)\ge p>0$ with probability at least $1-\delta/3$.
If $k\ge p^{-1}\log(3/\delta)$ then with probability at least $1-\delta$, $R_P(\hat r)\le R_P(r_H)+\eta+\varepsilon_m(k,\delta/3)$.
\end{proposition}

The procedure does not need to know $r_H$.

\section{Separation}\label{sec:separation}

\begin{definition}[Confusable family]\label{def:confusable}
$P_X$ on $\mathcal X$, $f_1,\dots,f_N:\mathcal X\to\mathcal Y$, $N\ge2$.
Task $P_\theta$ draws $X\sim P_X$ and sets $Y=f_\theta(X)$.
The family is $2\Delta$-separated if $P_X(f_\theta\ne f_{\theta'})\ge2\Delta$ for $\theta\ne\theta'$.
\end{definition}

\begin{lemma}\label{lem:confusable}
In a $2\Delta$-separated family every $r$ has $R_{P_\theta}(r)<\Delta$ for at most one $\theta$.
\end{lemma}

\begin{proposition}[Lower bound]\label{prop:lower-bound}
Every data-independent $r_H$ satisfies $\max_\theta R_{P_\theta}(r_H)\ge\Delta$ and $N^{-1}\sum_\theta\mathbb E\,R_{P_\theta}(r_H)\ge(1-1/N)\Delta$.
\end{proposition}

\begin{assumption}[LLM identifies the task]\label{ass:identifies}
There exist $\varepsilon_0\in[0,\Delta)$, $p>0$, $n_0$ such that for every $\theta$ and $n_1\ge n_0$, with probability at least $1-\delta/3$ over $S_1\sim P_\theta^{n_1}$, $Q(\{r:R_{P_\theta}(r)\le\varepsilon_0\}\mid S_1)\ge p$.
\end{assumption}

\begin{theorem}[Separation]\label{thm:separation}
Under Definition~\ref{def:confusable} and Assumption~\ref{ass:identifies}, let $n_1\ge n_0$ and $k\ge p^{-1}\log(3/\delta)$.
For every data-independent $r_H$ there exists $\theta$ such that with probability at least $1-\delta$, $R_{P_\theta}(r_H)-R_{P_\theta}(\hat r)\ge\Delta-\varepsilon_0-\varepsilon_m(k,\delta/3)$.
\end{theorem}
% TODO(check): if r_H is random, the task theta with R_{P_theta}(r_H) >= Delta depends on the realisation of r_H.
% The statement then needs theta chosen after r_H (conditionally on r_H), or the averaged form of Proposition prop:lower-bound.

The content of the theorem sits in Assumption~\ref{ass:identifies}.
That assumption is a statement about the LLM as a learner.

\section{Occam bound under the LLM prior}\label{sec:occam}

Let $\pi$ be a distribution on $\mathcal R$ fixed before seeing data, for example $\pi(r)=M(r\mid\text{description})$.
% TODO(check): M(r | description) is a distribution on R only if it includes the probability of ending the string after r.
% prior.py omits the end-of-text token (see the Experiments section).

\begin{proposition}[Occam bound under the LLM prior]\label{prop:occam}
With probability at least $1-\delta$, for all $r$, $R_P(r)\le\hat R_S(r)+\sqrt{(\log(1/\pi(r))+\log(1/\delta))/(2n)}$.
\end{proposition}

\begin{corollary}[Oracle inequality]\label{cor:oracle}
Let $B_\delta(r)$ be the square-root term and $\tilde r$ minimise $\hat R_S(r)+B_{\delta/2}(r)$.
For every fixed $r^\star$, with probability at least $1-\delta$, $R_P(\tilde r)\le R_P(r^\star)+2\sqrt{(\log(1/\pi(r^\star))+\log(2/\delta))/(2n)}$.
\end{corollary}

The sample complexity is of order $\varepsilon^{-2}\{\log(1/\pi(r^\star))+\log(1/\delta)\}$.
It depends on the description length of a good rubric under the LLM, not on the number of model parameters.
Taking $r^\star=r_H$ compares with the human prompt.

\begin{proposition}[PAC-Bayes and Gibbs posterior]\label{prop:pac-bayes}
With probability at least $1-\delta$, for all $\rho$, $\mathbb E_\rho R_P\le\mathbb E_\rho\hat R_S+\sqrt{(\mathrm{KL}(\rho\|\pi)+\log(2\sqrt n/\delta))/(2n)}$ \citep{maurer2004note}.
The Gibbs posterior $\rho_\lambda(r)\propto\pi(r)\exp\{-\lambda n\hat R_S(r)\}$ minimises $\mathbb E_\rho\hat R_S+\mathrm{KL}(\rho\|\pi)/(\lambda n)$.
\end{proposition}
% TODO(check): Theorem 5 of maurer2004note assumes n >= 8, which the statement omits.

It is an open question whether prompting the LLM with $S$ samples approximately from $\rho_\lambda$.
That would place the procedure inside generalised Bayesian inference \citep{bissiri2016general}.

\section{One-sided guarantees}\label{sec:one-sided}

Write $R^+_P(r)=\mathbb P(\hat Y=1\mid Y=0)$.
This mirrors trustable verification in \citet{balcan2025verifiers}.

\begin{proposition}[Realisable false-positive control]\label{prop:one-sided}
$S_2$ contains $m_-$ negatives.
Among candidates with no false positive on them, let $\hat r$ have the smallest empirical false-negative rate.
If some candidate has $R^+_P=0$, then with probability at least $1-\delta$, $R^+_P(\hat r)\le\log(k/\delta)/m_-$.
\end{proposition}


\section{Experiments}\label{sec:experiments}

\texttt{example1.py} illustrates Theorem~\ref{thm:separation} and Lemma~\ref{lem:confusable}.
It generates short item descriptions with independent uniform attributes in \texttt{data.py}, and three hidden rules on these attributes define a confusable family as in Definition~\ref{def:confusable}.
% TODO(check): the code estimates Delta as half the smallest pairwise disagreement on the held-out test inputs, not under P_X itself.
For each task it evaluates a few fixed handcrafted prompts, which are data-independent, and runs the rubric procedure of \texttt{rubric.py} with the language model as proposal kernel $Q$.
% TODO(check): R_P is replaced everywhere by the average loss on a fixed held-out test set, which is shared across tasks and seeds.
It reports the risk of $\hat r$ and of the best candidate, and estimates $p$ by the fraction of the $k$ candidates whose risk is at most $\varepsilon_0$.
It then counts the tasks on which each handcrafted prompt has risk below $\Delta$ and compares the largest gap in risk between the handcrafted prompt and $\hat r$ with $\Delta-\varepsilon_0-\varepsilon_m$.
% TODO(check): the code compares with Delta - eps0 - eps_m(k, delta), whereas Theorem thm:separation uses eps_m(k, delta/3).
% TODO(check): the code uses the mean of R(r_hat) over seeds, whereas the theorem is a high-probability statement for a single run.
% TODO(results): one bar plot per task comparing the handcrafted prompts with r_hat and the best candidate.

\texttt{example2.py} fixes one task.
Part (a) illustrates Proposition~\ref{prop:selection}.
It keeps $n_1$ and $k$ fixed, varies $m$, and compares the excess risk $R_P(\hat r)-\min_jR_P(r_j)$ over several seeds with $\varepsilon_m(k,\delta)$.
% TODO(results): excess risk against m, with eps_m(k, delta) on the same axes.
Part (b) illustrates Proposition~\ref{prop:occam} and Corollary~\ref{cor:oracle}.
It computes $\log(1/\pi(r))$ in \texttt{prior.py} as the negative log-probability of the rubric text under the same model, given only the task description.
% TODO(check): prior.py sums the log-probabilities of the tokens of r without an end-of-text token, so pi is not normalised over R and Proposition prop:occam does not apply as stated.
% TODO(check): without llama-cpp-python, prior.py uses log(1/pi(r)) = (c+1) log 2 + c log 128, which sums to at most one only if Sigma has at most 128 symbols, and generated rubrics may contain other characters.
For a pool of generated rubrics and handcrafted prompts it evaluates the bound on a fresh sample for several $n$, the fraction of rubrics whose risk lies below their bound, and the rubric $\tilde r$ with the smallest bound.
% TODO(check): the code selects r_tilde by minimising hat R_S + B_delta, whereas Corollary cor:oracle minimises hat R_S + B_{delta/2}.
% TODO(check): the code reports the fraction of covered rubrics, whereas Proposition prop:occam bounds the probability that all rubrics are covered.
% TODO(results): scatter of log(1/pi(r)) against risk, and R(r_tilde) with the bound at r_tilde against n.

\bibliographystyle{plainnat}
\bibliography{refs}

\appendix
\section{Proofs}\label{app:proofs}

\begin{proof}[Proof of Proposition~\ref{prop:selection}]
Condition on $S_1,r_1,\dots,r_k$.
Hoeffding and a union bound give $|\hat R_{S_2}(r_j)-R_P(r_j)|\le\sqrt{\log(2k/\delta)/(2m)}$ for all $j$.
Then
%
\[
R_P(\hat r)\le\hat R_{S_2}(\hat r)+\tfrac12\varepsilon_m\le\hat R_{S_2}(r_{j^\star})+\tfrac12\varepsilon_m\le R_P(r_{j^\star})+\varepsilon_m.
\qedhere
\]
\end{proof}
% TODO(check): j^star is not defined in the source. It is presumably an index attaining min_j R_P(r_j).

\begin{proof}[Proof of Corollary~\ref{cor:human-in-pool}]
\emph{To be written.}
% TODO(check): no proof supplied in the source.
\end{proof}

\begin{proof}[Proof of Proposition~\ref{prop:competitiveness}]
No candidate in $A_\eta$ has probability at most $(1-p)^k\le e^{-pk}\le\delta/3$.
Apply Proposition~\ref{prop:selection} at level $\delta/3$ and a union bound.
\end{proof}

\begin{proof}[Proof of Lemma~\ref{lem:confusable}]
Where $f_\theta\ne f_{\theta'}$ the labels differ, so $M(f_\theta(x)\mid r,x)+M(f_{\theta'}(x)\mid r,x)\le1$ and the two losses sum to at least one.
Hence $R_{P_\theta}(r)+R_{P_{\theta'}}(r)\ge2\Delta$.
\end{proof}

\begin{proof}[Proof of Proposition~\ref{prop:lower-bound}]
\emph{To be written.}
% TODO(check): no proof supplied in the source.
\end{proof}

\begin{proof}[Proof of Theorem~\ref{thm:separation}]
\emph{To be written.}
% TODO(check): no proof supplied in the source.
\end{proof}

\begin{proof}[Proof of Proposition~\ref{prop:occam}]
Hoeffding at level $\delta\pi(r)$ and a union bound.
\end{proof}

\begin{proof}[Proof of Corollary~\ref{cor:oracle}]
\emph{To be written.}
% TODO(check): no proof supplied in the source.
\end{proof}

\begin{proof}[Proof of Proposition~\ref{prop:pac-bayes}]
\emph{To be written.}
% TODO(check): no proof supplied in the source. The first part is cited from maurer2004note.
\end{proof}

\begin{proof}[Proof of Proposition~\ref{prop:one-sided}]
A candidate with $R^+_P>\varepsilon$ survives with probability at most $e^{-\varepsilon m_-}$.
Union over $k$ and set $\varepsilon=\log(k/\delta)/m_-$.
\end{proof}


\end{document}
